% Part: sets-functions-relations
% Chapter: sets
% Section: subsets

\documentclass[../../../include/open-logic-section]{subfiles}

\begin{document}

\olfileid{sfr}{set}{sub}
\olsection{उपसंच आणि घातसंच}

\begin{explain}
आपल्याला अनेकदा संचांची तुलना करावी लागेल. तुलना करण्याचा एक सरळ प्रकार
असा आहे: \emph{एका संचातील प्रत्येक वस्तू दुसऱ्या संचातही असते}.
ही स्थिती पुरेशी महत्त्वाची असल्यामुळे तिच्यासाठी आपण नवीन चिन्हपद्धतीची
ओळख करून घेऊ.
\end{explain}

\begin{defn}[उपसंच]
$A$ या संचातील प्रत्येक !!{element} हा $B$ चाही !!a{element} असेल,
तर $A$ हा $B$ चा \emph{उपसंच} आहे असे म्हणतो, आणि $A \subseteq B$ असे लिहितो.
$A$ हा $B$ चा उपसंच नसेल, तर $A \not\subseteq B$ असे लिहितो.
$A \subseteq B$ पण $A \neq B$ असेल, तर $A \subsetneq B$ असे लिहितो
आणि $A$ हा $B$ चा \emph{उचित उपसंच} आहे असे म्हणतो.
\end{defn}

\begin{ex}
प्रत्येक संच स्वतःचाच उपसंच असतो, आणि $\emptyset$ हा प्रत्येक संचाचा उपसंच असतो.
सम नैसर्गिक संख्यांचा संच हा नैसर्गिक संख्यांच्या संचाचा उपसंच आहे.
तसेच, $\{ a, b \} \subseteq \{ a, b, c \}$.
पण $\{ a, b, e
\}$ हा $\{ a, b, c \}$ चा उपसंच नाही.
\end{ex}

\begin{ex}
$2$ ही संख्या पूर्णांकांच्या संचाचा !!{element} आहे, तर सम संख्यांचा संच
हा पूर्णांकांच्या संचाचा उपसंच आहे. मात्र, एखादा संच दुसऱ्या संचाचा
!!a{element} आणि उपसंच \emph{दोन्हीही} असू शकतो. उदाहरणार्थ,
$\{0\} \in \{0, \{0\}\}$ आणि $\{0\} \subseteq \{0,
\{0\}\}$ देखील.
\end{ex}

विस्तारात्मकतेमुळे संचांच्या समानतेचा निकष मिळतो: $A = B$ तेव्हा आणि केवळ तेव्हाच,
जेव्हा $A$ मधील प्रत्येक !!{element} हा $B$ चाही !!a{element} असतो, आणि उलटही.
``उपसंच'' या व्याख्येत $A \subseteq B$ याचा अर्थ या निकषाचा नेमका पहिला
भाग आहे: $A$ मधील प्रत्येक !!{element} हा $B$ चा !!a{element} असतो.
अर्थात, $A$ आणि $B$ यांची अदलाबदल केली तरी ही व्याख्या लागू होते:
$B \subseteq A$ तेव्हा आणि केवळ तेव्हाच, जेव्हा $B$ मधील प्रत्येक !!{element}
हा $A$ चा !!a{element} असतो. हाच विस्तारात्मकतेतील ``आणि उलटही'' हा भाग होय.
दुसऱ्या शब्दांत, संच एकमेकांचे उपसंच असतील तेव्हा आणि केवळ तेव्हाच ते समान
असतात, हे विस्तारात्मकतेवरून निष्पन्न होते.

\begin{prop}
$A = B$ तेव्हा आणि केवळ तेव्हाच, जेव्हा $A \subseteq B$ आणि $B \subseteq A$ दोन्ही असतात.
\end{prop}

आणखी काही उपयुक्त चिन्हपद्धतींची ओळख करून घेण्याची ही योग्य वेळ आहे.
$A$ हा $B$ चा उपसंच केव्हा असतो याची व्याख्या करताना आपण
``$A$ मधील प्रत्येक !!{element} हा \dots'' असे म्हटले, आणि त्या
``$\dots$'' च्या जागी ``$B$ चा !!a{element} असतो'' असे घातले.
परंतु असा \emph{आकार} असणाऱ्या अभिव्यक्ती इतक्या नेहमी वापरल्या जातात की,
त्यांच्यासाठी आकारिक चिन्हपद्धती उपयोगी पडेल.

\begin{defn}\ollabel{forallxina}
$(\forall x \in A)\phi$ हे $\forall x(x \in A \lif
\phi)$ याचे संक्षिप्त रूप आहे. त्याचप्रमाणे, $(\exists x \in A)\phi$ हे
$\exists x(x
\in A \land \phi)$ याचे संक्षिप्त रूप आहे.
\end{defn}

या चिन्हपद्धतीचा उपयोग करून, $A \subseteq B$ तेव्हा आणि केवळ तेव्हाच, जेव्हा
$(\forall
x \in A)x \in B$, असे म्हणता येते.

आता आपण एका विशिष्ट प्रकारच्या संचाचा विचार करू: दिलेल्या संचाच्या सर्व
उपसंचांचा संच.

\begin{defn}[घातसंच]
$A$ या संचाच्या सर्व उपसंचांचा मिळून जो संच बनतो, त्याला
$A$ चा \emph{घातसंच} म्हणतात आणि तो $\Pow{A}$ असा लिहितात.
  \[
    \Pow{A} = \Setabs{B}{B \subseteq A}
  \]
\end{defn}

\begin{ex}
$\{ a, b, c \}$ चे सर्व शक्य उपसंच कोणते? ते असे आहेत:
$\emptyset$, $\{a \}$, $\{b\}$, $\{c\}$, $\{a, b\}$, $\{a, c\}$, $\{b,
c\}$, $\{a, b, c\}$. या सर्व उपसंचांचा संच म्हणजे
$\Pow{\{a,b,c\}}$:
\[
\Pow{\{ a, b, c \}} = \{\emptyset, \{a \}, \{b\}, \{c\}, \{a, b\},
\{b, c\}, \{a, c\}, \{a, b, c\}\}
\]
\end{ex}

\begin{prob}
$\{a, b, c, d\}$ चे सर्व उपसंच लिहा.
\end{prob}

\begin{prob}
$A$ मध्ये $n$ !!{element}s असतील, तर $\Pow{A}$ मध्ये $2^n$
!!{element}s असतात हे दाखवा.
\end{prob}

\end{document}
